% Created 2026-10-13 Tue 15:10
% Intended LaTeX compiler: lualatex
\documentclass[a4paper]{article}

\usepackage{amsmath}
\usepackage{fontspec}
\usepackage{graphicx}
\usepackage{longtable}
\usepackage{wrapfig}
\usepackage{rotating}
\usepackage[normalem]{ulem}
\usepackage{capt-of}
\usepackage{hyperref}
\usepackage[left=1in, right=1in, top=1in, bottom=1in]{geometry}
\usepackage{fontspec}
\usepackage{amsmath, amsfonts, amssymb}
\usepackage{amsthm}
\theoremstyle{definition}
\newtheorem{definition}{Definition}[section]
\newtheorem*{remark}{Remark}
\usepackage{pgf-umlsd}
\usepackage{tikz}
\usetikzlibrary{shapes,arrows,positioning,calc}
\usepackage{ebproof}
\newcommand{\msk}[1]{\mathsf{msk}_{#1}}
\newcommand{\mpk}[1]{\mathsf{mpk}_{#1}}
\newcommand{\usk}[2]{\mathsf{usk}_{#1,#2}}
\newcommand{\enc}[2]{\mathsf{enc}_{#1,#2}}
\newcommand{\sfusk}{\mathsf{usk}}
\newcommand{\sfenc}{\mathsf{enc}}
\newcommand{\iwrt}[1]{\,\Longrightarrow_{#1}\,}
\newcommand{\attr}[2]{#1_{#2}}
\newcommand{\powerset}[1]{\mathbb P(#1)}
\newcommand{\cA}{\mathcal{A}}
\newcommand{\cB}{\mathcal{B}}
\newcommand{\cC}{\mathcal{C}}
\newcommand{\cD}{\mathcal{D}}
\newcommand{\cH}{\mathcal{H}}
\newcommand{\cE}{\mathcal{E}}
\newcommand{\cK}{\mathcal{K}}
\newcommand{\cO}{\mathcal{O}}
\newcommand{\cP}{\mathcal{P}}
\newcommand{\cR}{\mathcal{R}}
\newcommand{\cS}{\mathcal{S}}
\newcommand{\cT}{\mathcal{T}}
\newcommand{\cU}{\mathcal{U}}
\newcommand{\cV}{\mathcal{V}}
\newcommand{\bN}{\mathbb{N}}
\newcommand{\true}{\mathsf{true}}
\newcommand{\false}{\mathsf{false}}
\newcommand{\True}{\mathsf{True}}
\newcommand{\False}{\mathsf{False}}
\date{}
\title{Covercrypt\\\medskip
\large An efficient ABE-like encryption scheme.}
\hypersetup{
 pdfauthor={},
 pdftitle={Covercrypt},
 pdfkeywords={},
 pdfsubject={},
 pdfcreator={},
 pdflang={English}}
\usepackage{biblatex}
\addbibresource{references.bib}
\begin{document}

\maketitle
\tableofcontents

\newpage
\section{Introduction}
\label{sec:orgd5c3498}

In large organizations, pieces of data created by one member are commonly shared
to other members.  However,  some pieces of data are sensitive  and must only be
shared to a specific subset of members. The process of accessing a piece of data
therefore becomes a three-step procedure  requiring members to: (1) authenticate
themselves  to  some  trusted  identity   management  system,  (2)  request  the
authorization to  access this  data from  some trusted  authorization management
system and  (3) request this data  from the storage system  which first verifies
the validity of the identity claim  delivered by the identity management system,
the  validity  of  the  authorization   claim  delivered  by  the  authorization
management system, and the compatibility of those claims.

However, authorization rules may be complex which may lead to operator errors or
bugs in  the authorization and  storage systems. Moreover actors  may misbehave,
which   may  question   our  trust   in   the  authorization   given  by   those
systems. Hopefully, there are ways to raise the need for trust in those systems.
One of them consists in  cryptographically enforcing authorization by encrypting
data under a key owned by the members  of the targeted only.  That way, the only
trusted parties are the targeted members  and the key distribution, which can be
performed in  a centralized Key  Management System (KMS).  But  if authorization
rules are complex, key distribution may  also become complex and hence victim of
the same kind of errors.

Attribute-Based Encryption (ABE) is a tool  that has been extensively studied to
solve exactly this  problem.  However, known ABE schemes are  not very efficient
and not post-quantum-secure which hinders  their adoption. Covercrypt solves the
same  problem  while  relaxing  some  formal constraints  in  order  to  achieve
efficiency in most practical use-cases  while offering both CCA and post-quantum
security, attribute-hiding ciphertexts and key tracing.

Covercrypt assumes  the existence  of a centralized  master authority  holding a
master secret key with  which all other keys are created:  the master public key
is used to encrypt data under its  access policy while user secret keys are also
associated to access policies and can  decrypt all Covercrypt ciphertexts with a
compatible access policy.  Covercrypt follows the KEM-DEM approach in which data
is symmetrically  encrypted under a key  which encapsulation is appended  to the
symmetric ciphertext. More precisely, Covercrypt is a actually Key-Encapsulation
Mechanism with  Access Control  (KEM-AC) --  as defined  in \ref{orgc8dd6ec} --
that can be  combined with a Data-Encryption Mechanism (DEM)  to obtain a Public
Key Encryption (PKE) scheme.  In the  rest of this document, we therefore prefer
the term Covercrypt \emph{encapsulation} to the term Covercrypt ciphertext. Moreover,
when necessary,  we use the  terms \emph{user  policy} and \emph{encapsulation  policy} to
denote  an  access  policy  associated  to   a  user  secret  and  a  Covercrypt
encapsulation respectively, and  use the term policy the context  makes it clear
whether we refer to a user or encapsulation policy.
\section{Preliminaries}
\label{sec:org1dda04b}

\subsection{Non-Interactive Key Exchange (NIKE)}
\label{sec:org594e973}

\begin{itemize}
\item Note taken on \textit{{[}2026-07-21 Tue 19:32] } \\
The key-homomorphic definition  from the ETSI document is  insufficient to allow
for an  easy computation of  \(\beta =  (s - \alpha  \cdot s_1) *  s_2^{-1}\). For
this, we really  need the notion of  field since secret keys  must be invertible
with respect  to the multiplicative  law.  Then, we  need the notion  that there
exists a  group generator \(\mathfrak  g\) for the public  keys such that  for any
pair of secret and public keys \((sk, pk)\), \(pk = sk \cdot \mathfrak g\).
\end{itemize}
\begin{definition}
A NIKE is  a cryptographic scheme allowing to privately  establish a session key
between  two  parties in  a  non-interactive  fashion  using the  following  two
primitives:
\begin{description}
\item[{keygen:}] \(\bot \rightarrow \cK_S \times \cK_P\)

This randomized algorithm  returns a new pair of secret  and public keys \((sk,
  pk)\).

\item[{session-key:}] \(\cK_P \times \cK_S \rightarrow \cK\)

Upon receiving a  public key \(pk \in  \cK_P\) and a secret key  \(sk \in \cK_S\),
returns the associated session  key \(K \in \cK\) such that  for any given pairs
of keys \((sk_1, pk_1)\)  and \((sk_2, pk_2)\), \(\mathrm{session\text{-}key}(pk_1,
  sk_2) = \mathrm{session\text{-}key}(pk_2, sk_1)\).
\end{description}
Additionally, a NIKE is said to be  \emph{key-homomorphic} if its public keys form an
Abelian group which scalars are its secret keys.
\end{definition}

\begin{remark}
One instantiation of a key-homomorphic NIKE  is based on an elliptic-curve group
\(G\) with prime order  \(p\) and group generator \(\mathfrak g\)  and a hash function
\(\cH : G \rightarrow \{0,1\}^\lambda\):

\begin{description}
\item[{keygen:}] returns \(sk  \overset{\mathdollar}{\longleftarrow} \cU(0, p-1)\) and
\(pk \gets sk \cdot \mathfrak g\).

\item[{session-key:}] upon receiving \(pk = sk \cdot \mathfrak g\) and \(sk'\), returns
\(K \gets \cH(sk \cdot sk' \cdot \mathfrak g)\).
\end{description}
\end{remark}
\subsection{Key-Encapsulation Mechanism (KEM)}
\label{sec:org1ff3c7f}

\begin{definition}
A KEM  is a cryptographic scheme  associated to the key  spaces \(\cK_D\), \(\cK_E\)
and \(\cK\), and encapsulation space \(\cE\), that exposes the following primitives:
\begin{description}
\item[{keygen:}] \(\bot \rightarrow \cK_D \times \cK_E\)

This  randomized   algorithm  returns   a  new   pair  of   decapsulation  and
encapsulation keys \((dk, ek)\).

\item[{encaps:}] \(\cK_P \rightarrow \cK \times \cE\)

This randomized algorithm  takes as input an encapsulation key  \(ek \in \cK_E\)
and returns a  new session key \(K  \in \cK\) and its  encapsulation \(\sfenc \in
  \cE\).

\item[{decaps:}] \(\cK_S \times \cE \rightarrow \cK\) or \(\bot\)

This  algorithm takes  as input  a  decapsulation key  \(dk \in  \cK_D\) and  an
encapsulation \(\sfenc \in  \cE\), and returns the encapsulated  session key \(K\)
if  and only  if  the  encapsulation was  created  for  the encapsulation  key
associated to \(dk\).
\end{description}
\end{definition}
\subsection{KEM with Access Control (KEM-AC)}
\label{sec:org4002894}

A Key-Encapsulation Mechanism with  Access Control (KEM-AC) allows encapsulating
a secret  with respect to an  \emph{encapsulation policy}, such that  user secret key
can open it  if an only if  its associated \emph{user policy} is  compatible with the
encapsulation policy with respect to some binary relation over the policy space.

\begin{definition}[Key-Encapsulation Mechanism with Access Control]
A KEM-AC  is defined over  a master-secret-key space  \(\cK_S\), master-public-key
space \(\cK_P\) and a user-secret-key space \(\cK_U\), a session-key space \(\cK\) and
an encapsulation space \(\cE\), and exposes the following operations:
\begin{description}
\item[{setup:}] \(\powerset{\cP_U \times \cP_E} \rightarrow \cK_S \times \cK_P\)

This randomized algorithm takes as input  a binary relation \(R\) over policies,
and returns a master secret key \(\msk{R}\) and a master public key \(\mpk{R}\).

\item[{keygen:}] \(\cK_S \times \cP_U \rightarrow \cK_U\)

This randomized  algorithm takes  as input  a master  secret key  \(\msk{R} \in
  \cK_S\) under  a binary relation  \(R \subset \cP_U  \times \cP_E\) and  a policy
\(p_u \in \cP_U\), and returns a user secret key \(\usk{p_u}{R}\).

\item[{encaps:}] \(\cK_P \times \cP_E \rightarrow \cK \times \cE\)

This randomized algorithm takes as input a master public key \(mpk_R \in \cK_P\)
under a  binary relation \(R \subset  \cP_U \times \cP_E\) and  an encapsulation
policy \(p_e  \in \cP_E\), and returns  both a session  key \(k \in \cK\)  and its
encapsulation \(\enc{p_e}{R} \in \cE\).

\item[{decaps:}] \(\cK_U \times \cE \rightarrow \cK\) or \(\bot\)

This algorithm  takes as input  a user  secret key \(\usk{p_u}{R_u}  \in \cK_U\)
associated to  an access policy \(p_u  \in \cP_U\) under a  binary relation \(R_u
  \subset \cP_U  \times \cP_E\),  and an  encapsulation \(\enc{p_e}{R_e}  \in \cE\)
associated to  a policy \(p_e \in  \cP_E\) under a binary  relation \(R_e \subset
  \cP_U \times \cP_E\).   It returns the encapsulated session key  \(k \in \cK\) if
\((p_u,p_e) \in R_e\), and \(\bot\) otherwise.
\end{description}
\label{orgc8dd6ec}
\end{definition}
\section{Subset-Cover KEM-AC (SC-KEM-AC)}
\label{sec:org58fca80}

\subsection{KEM-AC specification}
\label{sec:org58def2c}

\begin{itemize}
\item Note taken on \textit{{[}2026-10-13 Tue 14:55] } \\
\textbf{TODO}: List security claims somewhere.
\end{itemize}
Covercrypt is a  subset-cover algorithm: it implements a KEM-AC  with respect to
an  access structure  \(S\) defining  a  set of  rights \(\cU  = \{r_i\}_i\)  called
\emph{universe}, which policy space is  \(\powerset{\cU}\) and which binary relation is
the intersection test:  \(R_{\cU} = \{(p_u, p_e) \in \powerset{\cU}^2  : p_u \cap
p_e  \ne \emptyset  \}\).   The  complete reference  implementation  is given  in
Section \ref{sec:orgc79a45f} (note that the format of the keys is
dictated by the additional requirements listed in the following sections).
\subsection{Key and encapsulation life-cycle management}
\label{sec:orgb36d05b}

Just like any other keys, Covercrypt keys  can be compromised and lost. There is
therefore the need to renew them.  However,  in order to preserve the ability to
open older encapsulations with renewed keys, it is desirable for the key renewal
to be different from setting up a  new instance. The renewal procedure -- called
\emph{key rotation} -- is described in the present section.

Recall that  each user secret  key \(usk_p\) is associated  to a policy  \(p\). Upon
compromise of user secret key, a key  rotation with respect to its entire policy
must be  performed.  Due to  the particular  structure of Covercrypt  keys, this
process implies: (1) the rotation of all associated secrets in the master secret
key, (2) the  propagation of this rotation  to the master public  key to prevent
future encapsulations from being using  compromised secrets, (3) the propagation
of this rotation to all user secret keys which policy intersects the ones of the
compromised  key  to  allow  legitimate users  opening  the  new  encapsulations
generated with respect to the new secrets, and (4) optionally the refresh of all
encapsulations which associated policy intersects the one of the compromised key
to prevent  an eventual attacker from  opening them if it  still hasn't accessed
them.      The     entire     key-rotation    process     is     depicted     in
Figure \ref{rotation-flow}, and leverages the following primitives:
\begin{description}
\item[{usk-policy:}] \(\cK_U \rightarrow \powerset{\cU}\)

This function  takes as input a  Covercrypt user secret key  \(usk_p \in \cK_U\)
and returns its associated policy \(p \subset \cU\).

\item[{rotate:}] \(\cK_S \times \powerset{\cU} \rightarrow \cK_S \times \cK_P\)

This randomized  algorithm takes as input  a master secret key  \(\msk{\cU} \in
  \cK_S\) and an access policy \(p^* \subset \cU\), refreshes every secrets in this
given master secret key that are associated  to a right in this policy (\emph{i.e.}
each secret associated to a right in  \(s \cap p^*\)), and returns the refreshed
master secret key along with its associated master public key.

\item[{usk-refresh:}] \(\cK_S \times \cK_U \rightarrow \cK_U\) or \(\bot\)

This algorithm takes as  input a master secret key \(msk \in  \cK_S\) and a user
secret key \(usk_p  \in \cK_U\) associated to  a policy \(p \subset  \cU\). If the
integrity of the  given user secret key  can be verified by  the master secret
key, it returns a  new user secret key \(usk_p'\) associated  to the same access
policy containing all secrets that are contained in the given user secret key,
plus new secrets for each rights in  \(p\) for which newer secrets exists in the
master secret key.   If the integrity of  the given user secret  key cannot be
verified, \(\bot\) is returned.

\item[{enc-refresh:}] \(\cK_S \times \cK_P \times \cE \rightarrow \cE\) or \(\bot\)

This randomized algorithm takes as input  a master secret key \(msk \in \cK_S\),
a  master public  key \(mpk  \in \cK_S\)  and an  encapsulation \(enc_p  \in \cE\)
associated to a policy \(p \subset \cU\).  If the integrity of the master public
key  can be  verified, that  its revision  number matches  that of  the master
secret  key,  and  that the  encapsulation  can  be  open,  it returns  a  new
encapsulation \(enc_p'\) associated to the  same access policy. Otherwise \(\bot\)
is returned.
\end{description}

\tikzset{
  block/.style = {draw, fill=white, rectangle, minimum height=3em, minimum width=3em},
  tmp/.style  = {coordinate},
  sum/.style= {draw, fill=white, circle, node distance=2cm},
  input/.style = {coordinate},
  output/.style= {coordinate},
  pinstyle/.style = {pin edge={to-,thin,black}
  }
}

\begin{figure}
  \center
  \begin{tikzpicture}[auto, node distance=2cm,>=latex']
    \node [input] (usk) {};
    \node [block, right of=usk, node distance = 2.5cm] (pol) {usk-policy};
    \node [block, right of=pol, node distance = 2.5cm] (rot) {rotate};
    \node [input, below of=pol, node distance = 2.5cm] (enc) {};
    \node [block, below of=rot, node distance = 2.5cm, dashed] (erf) {enc-refresh};
    \node [input, above of=rot] (msk){};
    \node [output, below of=rot] (mout){};
    \node [output, below of=erf] (eout){};
    \node [block, right of=rot, node distance = 2.5cm] (urf) {usk-refresh};
    \node [input, above of=urf] (uski_old){};
    \node [output, below of=urf] (uski_new){};
    \draw [->] (rot) -- node{$msk'$} (urf);
    \draw [->] (rot) -- node{$\left(msk', mpk'\right)$} (mout);
    \draw [->] (erf) -- node{$enc_{p_j}$ or $\bot$} (eout);
    \draw [->] (uski_old) -- node{$usk_{p_i}^*$} (urf);
    \draw [->] (usk) -- node{$usk_{p^*}^*$} (pol);
    \draw [->] (enc) -- node{$enc_{p_j}^*$} (erf);
    \draw [->] (urf) -- node{$usk_{p_i}$ or $\bot$} (uski_new);
    \draw [->] (pol) -- node{$p^*$} (rot);
    \draw [->] (msk) -- node{$msk$} (rot);
  \end{tikzpicture}
  \caption{\label{rotation-flow}User-secret-key rotation flow, where untrusted objects are marked with an asterix.}
\end{figure}

\begin{remark}
Checking the  integrity of the user  secret key (resp.  encapsulation  key) upon
refresh is required  to prevent right escalation  (resp.  impersonification) via
modification of  this key, while  encapsulation policies cannot be  modified via
modification  of  the  encapsulation  due  to the  CCA  security  guaranteed  by
Covercrypt encapsulations.
\end{remark}

Additionally, to prevent impersonification attacks upon encapsulation, users can
always verify the integrity of the master public key using user secret keys, and
since this is necessary to prevent right-escalation attacks upon user-secret-key
refresh, we also provide a way to verify the integrity of user secret keys using
the master secret key:
\begin{description}
\item[{mpk-verify:}] \(\cK_U \times \cK_P \rightarrow \cB\)

Returns whether the  integrity of the given master public  key \(mpk \in \cK_P\)
can be verified using the given user secret key \(usk_p \in \cK_U\).

\item[{usk-verify:}] \(\cK_S \times \cK_P \rightarrow \cB\)

Returns whether the  integrity of the given user secret  key \(usk_p \in \cK_U\)
can be verified using the given master secret key \(msk \in \cK_S\).
\end{description}

Finally, a Covercrypt instantiation is  associated to a sequence of \emph{revisions}:
the setup initializes  an instance with \emph{revision number} 0,  and each execution
of the key rotation procedure increments this  number by one and returns the new
value of the master secret key with the new revision number.  Covercrypt objects
can only create objects with the same revision number which can be extracted via
one of the following functions:
\begin{description}
\item[{msk-revision:}] \(\cK_S \rightarrow \bN\)

Retrieves the  revision number from  the given  master secret key  \(\msk{} \in
  \cK_S\).

\item[{mpk-revision:}] \(\cK_P \rightarrow \bN\)

Retrieves the  revision number from  the given  master public key  \(\mpk{} \in
  \cK_P\).

\item[{usk-revision:}] \(\cK_U \rightarrow \bN\)

Retrieves the  revision number from  the given  master public key  \(\sfusk \in
  \cK_U\).

\item[{enc-revision:}] \(\cE \rightarrow \bN\)

Retrieves the revision number from the encapsulation \(\sfenc \in \cE\).
\end{description}

\begin{remark}
Downgrade attacks upon  encapsulation can be avoided by  verifying the integrity
of the master public key used, and that its revision number is no less than that
of the user secret key.
\end{remark}
\subsection{{\bfseries\sffamily TODO} Hybridization à la carte}
\label{sec:org31b1db6}

To accommodate different regulations and  right sensitivities, it is possible to
describe for each  right the kind of  encapsulation to generate, among  a set of
predefined encapsulation kinds.
\subsection{Live access-structure modifications}
\label{sec:orgd2db5b9}

Real-life access  structure may evolve  with time.  Covercrypt reflects  this by
exposing primitives  allowing to  modify the policy  universe without  having to
perform a new  setup. In particular, additive modifications of  the universe are
\emph{backward-compatible}:  old user  secret  keys and  encapsulations  need not  be
modified in order to open new compatible  encapsulations, or to be opened by new
compatible user secret keys. In other words, additive universe modifications can
be seen as the use of a right  that was present in the universe but unused until
then.
\begin{description}
\item[{add-rights:}] \(\cK_S \times \powerset{\cR} \rightarrow \cK_S \times \cK_P\)

This randomized algorithm  takes as input a master secret  key \(msk \in \cK_S\)
under a universe  \(\cU\) and a set  of new rights \(\{r_i\}_i  \subset \cR\), and
returns new  master secret and  public keys under  a universe \(\cU'  \gets \cU
  \cup \{r_i\}_i\).

\item[{del-rights:}] \(\cK_S \times \powerawr{\cR} \rightarrow \cK_S \times \cK_P\)

This randomized algorithm  takes as input a master secret  key \(msk \in \cK_S\)
under a universe  \(\cU\) and a set  of new rights \(\{r_i\}_i  \subset \cR\), and
returns new  master secret and  public keys with incremented  revision numbers
under a universe \(\cU' \gets \cU \setminus \{r_i\}_i\).
\end{description}

\begin{remark}
Additive universe modifications do not  generate new revisions while destructive
universe  modifications  generate  new  revisions since  downgrade  attacks  can
leverage  master public  keys  prior to  the  removal of  some  rights from  the
universe to trick users create encapsulations for compromised secrets.
\end{remark}
\subsection{SC-KEM-AC reference implementation}
\label{sec:orgc79a45f}

In this section, we give a reference implementation for the Covercrypt SC-KEM-AC
algorithm.   This implementation  is optimized  for  simplicity and  may not  be
optimized for speed. It presupposes the  existence of a key-homomorphic NIKE and
KEM schemes as described in the Section \ref{sec:org1dda04b}.
\subsubsection{Key format}
\label{sec:orge305681}

\subsubsection{SC-KEM-AC primitives}
\label{sec:orgac8ae39}

This implementation presuppose the existence of hash functions \(\mathtt{G, H, J,
U, T}\) which signatures are as follows:
\begin{itemize}
\item \(G : \{0,1\}^* \rightarrow \cK_S^\mathrm{nike}\)
\item \(T : (\cE^\mathrm{kem})^* \rightarrow \cT_T\)
\item \(U : (\{0,1\}^*)^* \rightarrow \cT_U\)
\item \(H : \cK^\mathrm{nike} \times \cK^{kem} \times \cK_P^\mathrm{nike} \times
  \cK_P^\mathrm{nike} \times \cT_T \rightarrow \{0,1\}^*\)
\item \(J : \{0,1\}^* \times \cK_P^\mathrm{nike} \times \cK_P^\mathrm{nike} \times
  \cT_U \rightarrow \{0,1\}^* \times \{0,1\}^*\)
\end{itemize}
Or, alternatively:
\begin{verbatim}
module type G = sig
  val hash : bytes -> Nike.sk
end

module type T = sig
  type t

  val hash : Kem.enc list -> t
end

module type U = sig
  type t

  val hash : bytes list -> t:T.t -> t
end

module type H = sig
  val hash : Nike.key -> Kem.key -> c1:Nike.pk -> c2:Nike.pk -> T.t -> bytes
end

module type J = sig
  type key
  type tag

  val hash : secret:bytes -> c1:Nike.pk -> c2:Nike.pk -> U.t -> key * tag
end
\end{verbatim}
Note that  the output of \texttt{T}  and \texttt{U} are  abstract and thus different  from any
other type, and that the \texttt{key} output of  \texttt{J} is the type of the key returned by
the  \texttt{encaps}  and  \texttt{decaps}  primitives.   With this  in  mind,  the  reference
implementation of the primitives of Covercrypt SC-KEM-AC is as follows:

\begin{description}
\item[{setup:}] \(\cU \rightarrow \cK_S \times \cK_P\)

Given the following master secret and public key types,

\begin{verbatim}
  rsk : (Nike.sk * Kem.dk) RightMM.t;
  seal : Sig.seal;
}

type enc = {
  v : nat;
  tag : tag;
  c1 : Nike.pk;
  c2 : Nike.pk;
  xenc : (Kem.enc * keylen sbytes) list;
}

let rpk_to_bytes rpk =
  RightMap.fold
    (fun right (pk, ek) bytes ->
      Right.to_bytes right :: Nike.pk_to_bytes pk :: Kem.ek_to_bytes ek
      :: bytes)
\end{verbatim}

the setup is a randomized algorithm that  takes as input a universe \(\cU\), and
returns  a   master  secret  key   \(\msk{R_\cU}\)  and  a  master   public  key
\(\mpk{R_\cU}\).

\begin{verbatim}
let encaps rng mpk policy =
  let pks =
    Policy.fold
      (fun right acc ->
        match RightMap.get mpk.rpk right with
        | Some (pk, ek) -> (pk, ek) :: acc
        | None -> invalid_arg "invalid policy")
      policy []
  in
  let secret = SBytes.gen rng keylen in
  let r = G.hash secret in
  let c1 = Nike.(r * mpk.h1) in
  let c2 = Nike.(r * mpk.h2) in
  let encs =
\end{verbatim}

Note the companion procedure \texttt{get\_mpk msk} given below can be used to the
generate master public keys from the master secret key after the setup.

\begin{verbatim}
let vk = Sig.get_vk msk.sk in
let rsk =
  Policy.fold
    (fun r ->
      match RightMM.get msk.rsk r with
      | None -> invalid_arg "invalid policy"
      | Some secrets -> RightMM.add r (NonEmptyList.head secrets))
    policy msk.rsk
in
let v = msk.v in
let seal = sign_usk ~sk:msk.sk ~v ~u1 ~u2 ~h1 ~h2 ~rsk in
{ v; u1; u2; h1; h2; vk; rsk; seal }
\end{verbatim}

\item[{keygen:}] \(\cK_S \times \cP_U \rightarrow \cK_U\)

\begin{verbatim}
  @@ List.map
       (fun (pk, ek) ->
         let k = Nike.session_key pk r in
         let k', enc = Kem.encaps rng ek in
         (k, k', enc))
       pks
in
let t = T.hash @@ List.map (fun (_, _, e) -> e) encs in
let xenc =
  List.map
    (fun (k, k', e) ->
      let h = H.hash k k' ~c1 ~c2 t in
      let f = SBytes.xor secret h in
      (e, f))
    encs
in
\end{verbatim}

\item[{encaps:}] \(\cK_P \times \cP_E \rightarrow \cK \times \cE\)

\begin{verbatim}
    let u = U.hash ~t @@ List.map Pair.snd xenc in
    J.hash ~secret ~c1 ~c2 u
  in
  (key, { v = mpk.v; tag; c1; c2; xenc })

let usk_secrets usk =
  interleave @@ List.of_seq
  @@ Seq.map NonEmptyList.to_list
  @@ RightMM.to_val_seq usk.rsk

let try_decaps enc p h1 h2 t u sk dk e f =
  match Kem.decaps dk e with
  | None -> None
  | Some k' ->
      let k = Nike.session_key p sk in
      let secret = SBytes.xor f @@ H.hash k k' ~c1:enc.c1 ~c2:enc.c2 t in
      let key, tag = J.hash ~secret ~c1:enc.c1 ~c2:enc.c2 u in
      if tag = enc.tag then
        let r = G.hash secret in
        let c1 = Nike.(r * h1) in
        let c2 = Nike.(r * h2) in
        if enc.c1 = c1 && enc.c2 = c2 then Some key else None
      else None

let decaps usk enc =
  let p1 = Nike.(usk.u1 * enc.c1) in
  let p2 = Nike.(usk.u2 * enc.c2) in
  let p = Nike.Point.(p1 + p2) in
  let t = T.hash @@ List.map (fun (e, _) -> e) enc.xenc in
  let u = U.hash ~t @@ List.map Pair.snd enc.xenc in
  List.fold_left
    (fun key (sk, dk) ->
\end{verbatim}

\item[{decaps:}] \(\cK_U \times \cE \rightarrow \cK\) or \(\bot\)

\begin{verbatim}
let p = Nike.Point.(p1 + p2) in
let t = T.hash @@ List.map (fun (e, _) -> e) enc.xenc in
let u = U.hash ~t @@ List.map Pair.snd enc.xenc in
List.fold_left
  (fun acc (r, sk, dk) ->
    List.fold_left
      (fun acc (e, f) ->
        match (try_decaps enc p h1 h2 t u sk dk e f, acc) with
        | None, acc -> acc
        | Some key, None -> Some (key, Policy.add r Policy.empty)
        | Some key, Some (key', policy) ->
            if key = key' then Some (key, Policy.add r policy)
            else invalid_arg "invalid encapsulation")
      acc enc.xenc)
\end{verbatim}

Where \texttt{try\_decaps} attempts opening a right encapsulation as follows:

\begin{verbatim}
let msk_secrets (msk : msk) =
  (* TODO: once a key in a history has opened the encapsulation, other keys
       from the same history need to attempt opening it again. *)
  interleave @@ List.of_seq
  @@ Seq.map (fun (r, secrets) ->
      List.map (fun (sk, dk) -> (r, sk, dk)) (NonEmptyList.to_list secrets))
  @@ RightMM.to_seq msk.rsk

let enc_policy msk enc =
  let p1 = Nike.(enc.c1) in
  let p2 = Nike.(Scalar.((msk.s - msk.s1) / msk.s2) * enc.c2) in
  let h1 = Nike.get_pk msk.s1 in
\end{verbatim}
\end{description}
\subsubsection{Key-rotation primitives}
\label{sec:orga79161b}

\begin{description}
\item[{usk-policy:}] \(\cK_U \rightarrow \powerset{\cU}\)

Extracting the policy out of a user secret key simply resorts to retrieving
the rights for which it holds secrets:

\begin{verbatim}
else if msk.v != mpk.v then invalid_arg "outdated master public key"
else
\end{verbatim}

\item[{rotate:}] \(\cK_S \times \poswerset{\cU} \rightarrow \cK_S \times \cK_P\)

\begin{verbatim}
    |> Option.map (fun (old_key, policy) ->
        let new_key, enc = encaps rng mpk policy in
        (~old_key, ~new_key, enc))

let msk_revision (msk : msk) = msk.v
let mpk_revision (mpk : mpk) = mpk.v
let usk_revision (usk : usk) = usk.v
let enc_revision (enc : enc) = enc.v
\end{verbatim}

\item[{usk-refresh:}] \(\cK_S \times \cK_U \rightarrow \cK_U\) or \(\bot\)

\begin{verbatim}
let usk_policy usk =
  RightMM.fold (fun r _ -> Policy.add r) usk.rsk Policy.empty

let rotate rng (msk : msk) policy =
  let v = Int32.succ msk.v in
  let rsk =
    Policy.fold
      (fun r -> RightMM.add r (Nike.gen_sk rng, Kem.gen_dk rng))
      policy msk.rsk
  in
  { msk with v; rsk }

let usk_refresh (msk : msk) usk keep_old_secrets =
  if not (is_valid_usk msk usk) then invalid_arg "invalid user secret key"
  else
    let v = msk.v in
    let rsk =
      RightMM.fold
        (fun r secrets rsk ->
          match RightMM.get msk.rsk r with
          | None -> rsk
          | Some master_secrets ->
              if keep_old_secrets then
                let newer_secrets =
                  NonEmptyList.take_while
                    (fun s -> s != NonEmptyList.head secrets)
                    master_secrets
\end{verbatim}

\item[{enc-refresh:}] \(\cK_S \times \cK_P \times \cE \rightarrow (\cK \times \cE)\)
or \(\bot\)

\begin{verbatim}
List.fold_left
  (fun acc (e, f) ->
    match (try_decaps enc p h1 h2 t u sk dk e f, acc) with
    | None, acc -> acc
    | Some key, None -> Some (key, Policy.add r Policy.empty)
    | Some key, Some (key', policy) ->
        if key = key' then Some (key, Policy.add r policy)
        else invalid_arg "invalid encapsulation")
  acc enc.xenc)
\end{verbatim}

Where \texttt{enc\_policy} allows retrieving the policy of an encapsulation along with
the encapsulated key:

\begin{verbatim}
        enc.xenc key)
    None
  @@ usk_secrets usk

let msk_secrets (msk : msk) =
  (* TODO: once a key in a history has opened the encapsulation, other keys
       from the same history need to attempt opening it again. *)
  interleave @@ List.of_seq
  @@ Seq.map (fun (r, secrets) ->
      List.map (fun (sk, dk) -> (r, sk, dk)) (NonEmptyList.to_list secrets))
  @@ RightMM.to_seq msk.rsk

let enc_policy msk enc =
  let p1 = Nike.(enc.c1) in
  let p2 = Nike.(Scalar.((msk.s - msk.s1) / msk.s2) * enc.c2) in
  let h1 = Nike.get_pk msk.s1 in
  let h2 = Nike.get_pk msk.s2 in
  let p = Nike.Point.(p1 + p2) in
  let t = T.hash @@ List.map (fun (e, _) -> e) enc.xenc in
  let u = U.hash ~t @@ List.map (fun (_, f) -> f) enc.xenc in
  List.fold_left
\end{verbatim}
\end{description}
\subsubsection{Live-universe-modification primitive}
\label{sec:org70fa593}
\section{Boolean KEM-AC}
\label{sec:org7501b11}

In order to  make expressing policies more ergonomic, we  introduce KEM-AC which
policy space  is a  set of  logical propositions  that is  closed under  the two
logical  operators \emph{and}  (\(\land\)) and  \emph{or} (\(\lor\)).   Then, we  show how  to
implement  this  boolean KEM-AC  on  top  of  the  set-based KEM-AC  defined  in
Section \ref{sec:org58fca80}  via compiling  boolean  policies  to
set-based policies.
\subsection{Syntax and semantics}
\label{sec:org7c3e9c9}

The policy language has a very simple syntax which specification is given below,
and which semantics is defined thereafter.
\phantomsection
\label{org95a3b6d}
\begin{verbatim}
<policy>    ::= <attribute> | <policy> <operator> <policy>
<operator>  ::= "&&" | "||"
<attribute> ::= "*" | <dimension>::<name>
<dimension> ::= <name>
<name>      ::= <word> \ { "&&" | "||" | <word>::<word> }
<word>      ::= char+
\end{verbatim}

\begin{definition}[Universe]
In the boolean KEM-AC setting, the universe  \(\Omega\) is a set of atomic logical
propositions called \emph{attributes}.
\end{definition}

\begin{definition}[Policies]
The  space of  policies  \(\cP_\Omega\)  associated to  a  universe of  attributes
\(\Omega\) is the closure of \(\Omega\) with respect to both logical operators \emph{and}
(\(\land\))  and \emph{or}  (\(\lor\)), noted  \(\cP_\Omega =  \left< \Omega,  \land, \lor
\right>\). In other words,  a policy is a function of sets  of true attributes in
\(\Omega\)   to   the   set   of  boolean   values   \(\cB\):   \(\cP_\Omega   \simeq
\powerset{\Omega}^\cB\). Given  a set of  true attributes \(X \subset  \Omega\), we
therefore use \(p(X)\)  to denote the value  taken by a policy  \(p \in \cP_\Omega\)
when exactly those attributes are set to true.
\end{definition}

\begin{remark}
\(\cP_\Omega\)  has two  special policies,  noted \(\True\)  and \(\False\),  that are
respectively always true and false.
\end{remark}

In order to structure the set of expressible policies, the values attributes can
take  can be  constrained  by the  definition of  an  access structure  grouping
logically-related  attributes  together  by either  anarchical  or  hierarchical
dimensions.

\begin{definition}[Dimension]
A  dimension \(D\)  of a  universe of  attributes \(\Omega\)  is a  constraint on  a
non-empty set of attributes from \(\Omega\) requiring that at least one of them be
true:
\begin{equation}
  \tag{dimension}
  \bigvee_{a \in D} a = \True
\end{equation}
The set  of dimensions of \(\Omega\), noted \(\cD_\Omega\),  is therefore the
set of its non-empty parts:  \(\cD_\Omega = \powerset{\Omega}^*\).
\end{definition}

\begin{definition}[Hierarchy]
A dimension \(D \in \cD_\Omega\) defined by  an access structure \(S\) is said to be
a hierarchy if it is associated  a bijective function \(\iota \in D^{[1,|D|]}\) in
\(S\).  In that  case, it additionally requires that the  logical implication be a
strict total order over \(D\):
\begin{equation}
  \tag{hierarchy}
  \forall (a, a') \in D^2, a  \implies a' \iff \iota(a) > \iota(a')
\end{equation}
\end{definition}

\begin{definition}[Anarchy]
Otherwise,  a  dimension  \(D \in  \cD_\Omega\)  is  said  to  be an  anarchy  and
additionally requires its attributes to be mutually exclusive:
\begin{equation}
  \tag{anarchy}
  \forall (a, a') \in D^2,
  \left( a = a' \right) \lor \left( a \implies \neg a' \right)
\end{equation}
\end{definition}

Not all combinations of dimensions can be defined by a given access structure.

\begin{definition}[Access Structure]
In the boolean KEM-AC setting, an access structure is a set of mutually disjoint
dimensions: \[ \cS = \{S \in \powerset{\cD} \mid \forall (D_i, D_j) \in S^2, D_i
\cap D_j = \emptyset \}\]
\end{definition}

\begin{definition}[Set of Expressible Postulates]
The set of expressible postulates \(\cU_S \subset \powerset{\Omega}\) described by
an access  structure \(S\) is  the set  of parts of  \(\Omega\) that are  valid with
respect to each dimension in \(S\).
\end{definition}

\begin{definition}[Valid Policy]
A policy \(p\) is said to be valid  with respect to an access structure \(S\), noted
\(p \in \cV_S\) if and only if  there exists a postulate in \(\cU_S\) that validates
\(p\): \[ p \in \cV_S \iff \exists X \in \cU_S : p(X) = \true \]
\end{definition}
\subsection{KEM-AC}
\label{sec:orga1fb084}

In the boolean KEM-AC setting, a  user secret key \(\usk{p_u}{S}\) associated to a
policy \(p_u \in \cP_\cU\) can open an encapsulation \(\enc{p_e}{S'}\) associated to
a policy \(p_e \in \cP_\Omega\) if and  only if the conjunction of \(p_u\) and \(p_e\)
is valid with respect to both \(S\) and  \(S'\): \[ (p_u, p_e) \in R \iff (p_u \land
p_e) \in (\cV_S \cap \cV_{S'})\]

\begin{remark}
It is  important to distinguish the  access structure in which  an encapsulation
was  created from  the one  in which  a user  secret key  was created  since the
universe of expressible postulates can be modified  in such a way that an access
policy that is valid w.r.t.  a first access structure \(S\) becomes invalid w.r.t.
the new access structure \(S'\).
\end{remark}
\subsubsection{Examples}
\label{sec:org6a8b19c}

\begin{enumerate}
\item Without access structure modification
\label{sec:org0a167df}

Let \(S = \{A, (H, \iota)\}\), where \(A = \{ \attr{a}{1}, \attr{a}{2}, \attr{a}{3}
\}\) is an  anarchical dimension, and \(H  = \{ \attr{1}{h}, \attr{2}{h}  \}\) is a
hierarchical dimension with \(\iota(\attr{2}{h}) > \iota(\attr{1}{h})\).

\begin{itemize}
\item Let  \(p_u  = \attr{a}{1}  \land  \attr{2}{h}\)  and  \(p_e =  \attr{a}{1}  \land
  \attr{2}{h}\). In this  case, a user secret key \(\usk{p_u}{\cU}\)  \emph{can} open an
encapsulation  \(\enc{p_e}{\cU}\)  since  \(p_u  \land p_e  =  \attr{a}{1}  \land
  \attr{2}{h} \ne \False\).

\item Let \(p_u =  \attr{a}{1} = \attr{a}{1} \land \attr{1}{h}\)  using the definition
of a hierarchical dimension. Let \(p_e = \attr{2}{h} = \attr{2}{h} \land \left(
  \lor_{i=1}^{3}   \attr{a}{i}  \right)   =   (\attr{2}{h}  \land   \attr{a}{1})
  \lor_{i=2}^{3}  \left(  \attr{2}{h}  \land   \attr{a}{i}  \right)\)  using  the
definition of an anarchical dimension.

In  this   case,  a  user   secret  key  \(\usk{p_u}{\cU}\)  \emph{cannot}   open  an
encapsulation   \(\enc{p_e}{\cU}\)    since   \(\neg    (\attr{1}{h}   \iwrt{\cU}
  \attr{2}{h})\) by definition of a hierarchical dimension.
\end{itemize}
\item With expressible universe modification
\label{sec:org0c06f0e}

\begin{enumerate}
\item Additive transformation
\label{sec:org58e6592}

Let \(\cU' = \cU  \cup D_3\), where \(D_3 = \{  \attr{3}{1}, \attr{3}{2} \} \subset
\cA_\Omega\) is an anarchical dimension.
\begin{itemize}
\item Let  \(p_u  = \attr{1}{1}  \land  \attr{2}{2}\)  and  \(p_e =  \attr{1}{1}  \land
  \attr{2}{2}  \land  \attr{3}{2}\).  Then  a  user  secret key  \(\usk{p_u}{\cU}\)
\emph{cannot} decrypt  \(\enc{p_e}{\cU'}\) since  \(p_e =_{\cU} \False\)  and \(\neg(p_u
  \iwrt{\cU} \False)\).

\(p_u   =_{\cU'}  \attr{1}{1}   \land  \attr{2}{2}   \land  (\attr{3}{1}   \lor
  \attr{3}{2})\)

\item Let  \(p_u  = \attr{1}{1}  \land  \attr{2}{2}\)  and  \(p_e =  \attr{1}{1}  \land
  \attr{2}{2}  \land  \attr{3}{2}\).  Then  a user  secret  key  \(\usk{p_u}{\cU}\)
\emph{cannot} decrypt \(\enc{p_e}{\cU'}\) since
\end{itemize}
\item Destructive transformation
\label{sec:org53b21db}
\end{enumerate}
\end{enumerate}
\subsection{{\bfseries\sffamily TODO} Live universe modifications}
\label{sec:org1d3d2bc}

Evolving  the access  structure  underlying  a given  revision  of a  Covercrypt
instance can be done through one of those functions:
\begin{description}
\item[{add-hierarchy:}] \(\cK_S \times \powerset{\cR} \rightarrow
  \cK_S \times \cK_P\) or \(\bot\)
\item[{del-hierarchy:}] \(\cK_S \times \powerset{\cR} \rightarrow
  \cK_S \times \cK_P\) or \(\bot\)

This randomized algorithm  takes as input a master secret  key \(msk \in \cK_S\)
with an access structure \(S\) and a  set of attributes \(D \subset \cR\). If this
new hierarchy is valid with respect to \(S\), it returns a new master secret key
and a new master public key which structure contains (resp.  does not contain)
\(H\). Otherwise \(\bot\) is returned.

\item[{add-anarchy:}] \(\cK_S \times \cA \rightarrow \cK_S
  \times \cK_P\) or \(\bot\)
\item[{del-anarchy:}] \(\cK_S \times \cA \rightarrow \cK_S
  \times \cK_P\) or \(\bot\)

This randomized algorithm  takes as input a master secret  key \(msk \in \cK_S\)
with an access structure \(S\) and an  anarchical dimension \(A \in \cA\). If this
new anarchy is valid  with respect to \(S\), it returns a  new master secret key
and a new master public key which structure contains (resp.  does not contain)
\(A\). Otherwise \(\bot\) is returned.

\item[{add-attribute:}] \(\cK_S \times \cA \times \cU
  \rightarrow \cK_S \times \cK_P\) or \(\bot\)
\item[{del-attribute:}] \(\cK_S \times \cA \times \cU
  \rightarrow \cK_S \times \cK_P\) or \(\bot\)

This randomized algorithm  takes as input a master secret  key \(msk \in \cK_S\)
with  an access  structure \(S\),  an anarchical  dimension \(D  \in \cA\)  and an
attribute \(p \in \cU\). If this dimension is  a member of \(S\), it returns a new
master secret  key and  a new master  public key with  an access  structure in
which this dimension contains (resp. does not contain) this attribute.

\item[{add-hierarchical-attribute:}] \(\cK_S \times \cH \times \cU \rightarrow \cK_S
  \times \cK_P\) or \(\bot\)
\item[{del-hierarchical-attribute:}] \(\cK_S \times \cH \times \cU \rightarrow \cK_S
  \times \cK_P\) or \(\bot\)

This randomized algorithm  takes as input a master secret  key \(msk \in \cK_S\)
with an  access structure  \(S\), a  hierarchical dimension \(D  \in \cH\)  and an
attribute \(p \in \cU\). If this dimension is  a member of \(S\), it returns a new
master secret  key and  a new master  public key with  an access  structure in
which this dimension contains (resp. does not contain) this attribute.
\end{description}
\subsection{{\bfseries\sffamily TODO} Hybridization à la carte}
\label{sec:orgba73d5b}
\section{{\bfseries\sffamily TODO} Optional key tracing}
\label{sec:orgfb3bfd1}
\section{{\bfseries\sffamily TODO} Cooperating Covercrypt authorities}
\label{sec:org72bef40}

Covercrypt  suppose  the  existence  of   an  central  authority  in  charge  of
maintaining the master secret key, and  generating the master public key and all
user secret keys. However, in some use  cases, it is desirable to distribute the
power over  multiple authorities: users managed  by a given authority  should be
able to create encapsulations for attributes defined by a different authority.

Doing so requires  (1) generating fragments of a master  public key by selecting
the subset of the rights from the original key that is targeted for sharing, and
(2) combining master public key fragments together to create a new master public
key:
\begin{description}
\item[{\texttt{select:}}] \texttt{mpk -> policy -> mpk}
\item[{\texttt{combine:}}] \texttt{mpk -> mpk -> mpk}
\end{description}
Doing so requires a broader acception of  what a master public key is since such
key may not contain all rights hold by its associated master secret key, and may
not even be associated to a single master secret key anymore.

\newpage
\appendix
\end{document}
